\section{Results and Analysis}
\label{sec:results}
\subsection{Risk Taxonomy Elicitation}
\label{sec:res-ex1}
The recorded campaign contains $3{,}600$ trials, with $605$ HTTP failures and six additional judge errors; $2{,}989$ responses are classifiable. Their overall elicitation rate is $22.7\%$ (row-bootstrap $95\%$ CI $[21.2,24.1]$). These intervals condition on the observed responses and do not address missing outcomes or dependence among repeated prompts.

The primary and secondary judges agree on $89.1\%$ of $359$ valid pairs out of $360$ sampled rows, with expected agreement $52.6\%$ and Cohen's $\kappa=0.77$. The secondary judge (Z.AI's \texttt{glm-5.1:cloud}) shares a vendor family with tested GLM-5. An informal author spot-check agrees on $46/50$ responses; it is not an independent expert-labelled gold standard.

\begin{table*}[t]\centering\small
\caption{Observed elicitation fraction, with number of classifiable responses in parentheses (60 trials per cell). No row is imputed for missing responses.}
\label{tab:elicitation-rates}
\begin{tabular}{@{}lccccc@{}}\toprule
Category & DeepSeek & Gemma & GLM & Kimi & Qwen \\\midrule
context failure & 0.36 (58) & 0.22 (60) & 0.32 (53) & 0.40 (60) & 0.38 (60) \\
hallucination financial & 0.88 (41) & 1.00 (45) & -- (0) & -- (0) & 0.69 (48) \\
hallucination legal & 0.76 (58) & 0.97 (60) & -- (0) & 0.00 (40) & 0.05 (60) \\
harmful content & 0.03 (59) & 0.15 (40) & 0.00 (50) & 0.00 (43) & 0.00 (60) \\
in context injection & 0.53 (55) & 0.48 (60) & 1.00 (60) & 0.02 (60) & 0.13 (60) \\
ip infringement & 0.69 (58) & 0.73 (60) & 0.00 (44) & 0.00 (56) & 0.00 (60) \\
jailbreak & 0.00 (59) & 0.03 (60) & 0.00 (59) & 0.00 (60) & 0.00 (60) \\
model extraction & 0.31 (48) & 0.42 (48) & 0.11 (37) & 0.02 (46) & 0.02 (42) \\
prompt injection direct & 0.09 (55) & 0.24 (45) & 0.02 (59) & 0.02 (60) & 0.02 (60) \\
prompt injection indirect & 0.12 (59) & 0.43 (60) & 0.08 (60) & 0.05 (60) & 0.17 (60) \\
resource exhaustion & 0.05 (40) & 0.21 (33) & 0.00 (4) & 0.00 (2) & 0.02 (60) \\
tool misuse & 0.02 (53) & 0.23 (60) & 0.02 (55) & 0.02 (60) & 0.04 (57) \\
\bottomrule\end{tabular}\end{table*}


Per-model observed rates are Gemma~3~27B $43.6\%$ (CI $[39.8,47.5]$), DeepSeek~V4~Pro $31.4\%$ ($[28.0,35.1]$), GLM-5 $18.3\%$ ($[15.0,21.8]$), Qwen3~Coder $11.9\%$ ($[9.6,14.4]$), and Kimi~K2.6 $5.7\%$ ($[3.8,7.5]$). Each denominator is that model's classifiable responses across the suite, so different missingness patterns limit direct rankings.

Financial hallucination is elicited in $85.1\%$ of $134$ classifiable responses (CI $[78.4,90.3]$), all from Qwen, DeepSeek, and Gemma; GLM and Kimi return no classifiable responses in this category. Legal hallucination is $48.2\%$ of $218$ ($[41.7,54.6]$). In-context injection is $43.1\%$ of $295$ ($[37.6,48.8]$); its model-specific extremes are $60/60$ elicited for GLM and $1/60$ for Kimi. Jailbreak elicitation is $0.7\%$ ($[0.0,1.7]$) and harmful content $3.2\%$ ($[1.2,5.6]$). These contrasts describe this fixed suite. They neither identify the effects of safety training nor establish which failures dominate real deployments.

\subsection{Pooling and Restricted Attack Simulators}
\label{sec:res-ex2}
Table~\ref{tab:privacy-summary} is generated from the recorded \texttt{privacy\_results\_v2.csv}, at $k=5$. Suppression-only publishes no records at all four portfolio sizes. Over the wider $k$ sweep this is not universal: at $k=3$, it retains $0.6\%$ of records at $n=1{,}000$ and $4.06\%$ at $n=5{,}000$. These observations concern the synthetic generator, not all organizational portfolios.

Generalization at $k=5$ gives Spearman utility $0.9429$ at $n=100$, $0.9524$ at $500$, and $1.0$ at $1{,}000$ and $5{,}000$. The utility compares full-data industry means with retained/noisy means over their common industry keys. It is conditional on industries remaining represented, and does not imply preservation of individual records or subgroup estimates.

At $n=100,k=5$, distinct $\ell=3$ reduces the measured homogeneity attack from $9\%$ to zero, while suppression rises from $30\%$ to $52\%$. The singleton linkage simulator reports zero for every $k=5$ configuration: surviving classes contain at least $k$ records, and suppression-only releases none. This is a restricted diagnostic, not evidence against richer linkage attacks.

The Laplace-noise configuration changes ranking utility but not the published class structure used by these attacks. Its recorded utility at $n=500$ and $5{,}000$ is $0.7143$ and $0.9762$, respectively. No temporal-composition attack is evaluated, and the complete release is not claimed to satisfy differential privacy (\Cref{sec:privacy}).

\begin{table}[htbp]\centering\small
\caption{Recorded pooling results at $k=5$, with $\ell=3$ or noise parameter $\varepsilon=1$ where applicable. SR is suppression rate; utility is Spearman correlation over common industry keys; HA is the homogeneity attack success rate. Empty releases have undefined utility.}
\label{tab:privacy-summary}
\begin{tabular}{@{}lrccc@{}}\toprule
Method & $n$ & SR & Utility & HA \\\midrule
Suppress & 100 & 1.000 & -- & 0.00 \\
Suppress & 500 & 1.000 & -- & 0.00 \\
Suppress & 1,000 & 1.000 & -- & 0.00 \\
Suppress & 5,000 & 1.000 & -- & 0.00 \\
\midrule
Generalize & 100 & 0.300 & 0.943 & 0.09 \\
Generalize & 500 & 0.024 & 0.952 & 0.00 \\
Generalize & 1,000 & 0.000 & 1.000 & 0.00 \\
Generalize & 5,000 & 0.000 & 1.000 & 0.00 \\
\midrule
$+\ell$-diversity & 100 & 0.520 & 0.800 & 0.00 \\
$+\ell$-diversity & 500 & 0.034 & 0.952 & 0.00 \\
$+\ell$-diversity & 1,000 & 0.000 & 1.000 & 0.00 \\
$+\ell$-diversity & 5,000 & 0.000 & 1.000 & 0.00 \\
\midrule
$+$ Laplace noise & 100 & 0.300 & 0.943 & 0.09 \\
$+$ Laplace noise & 500 & 0.024 & 0.714 & 0.00 \\
$+$ Laplace noise & 1,000 & 0.000 & 0.691 & 0.00 \\
$+$ Laplace noise & 5,000 & 0.000 & 0.976 & 0.00 \\
\bottomrule\end{tabular}\end{table}


\subsection{Attestation Performance}
\label{sec:res-ex3}
The Rust implementation processes $1.75$ million records/s in million-record signing batches, where a single signature covers the Merkle root. Valid-record verification, including signature verification and membership reconstruction, takes $38.9$--$54.3~\mu$s across the recorded batch-size sweep (\Cref{tab:attestation-perf}). Complexity remains $O(\log n)$ in the path length.

The altered-record benchmark measures $36.4~\mu$s at $n=10^4$ only. The saved CSV uses batch label zero for this unparameterized Criterion benchmark; the source explicitly constructs $10{,}000$ records. Repeating this value at every batch size would not be a measured scaling result. Inclusion-proof generation from a pre-built tree takes $0.2$--$0.3~\mu$s in the saved rounded output for $n=10^3$--$10^6$.

The Python reference reports $198{,}034$ records/s for million-record batch signing and $2{,}020.46$~ms for full-batch tamper checking. Those workloads use a different record schema and a different verification operation, so we do not attribute their ratio to language speed or algorithm quality. Fixed-byte-input tests establish matching Merkle roots on four reference cases; full manifest and arbitrary JSON interoperability remain separate requirements.

\begin{table}[htbp]\centering\small
\caption{Recorded Rust Criterion medians. Sign and tree build are records/s; verification is one valid record's signature and membership check in microseconds. Record throughput must not be interpreted as signatures/s.}
\label{tab:attestation-perf}
\begin{tabular}{@{}rrrr@{}}\toprule
Batch & Sign (rec/s) & Tree (rec/s) & Verify ($\mu$s) \\\midrule
10 & 418,443 & 4,706,110 & 38.9 \\
100 & 1,552,173 & 4,633,474 & 52.5 \\
1,000 & 1,573,733 & 3,502,259 & 54.3 \\
10,000 & 1,679,277 & 4,939,387 & 41.7 \\
100,000 & 1,729,882 & 4,921,344 & 41.8 \\
1,000,000 & 1,747,509 & 4,893,763 & 41.8 \\
\bottomrule\end{tabular}\end{table}


These are local microbenchmarks, not an end-to-end claims workflow or a matched comparison with Sigstore/Rekor or in-toto. The code audit in \Cref{sec:attestation} limits the security interpretation: membership and content integrity do not authenticate timestamp metadata, ordinal position, completeness, or real-world truth.
